\documentclass[sigconf]{acmart}
\usepackage{algpseudocode}
\usepackage{float}
\usepackage{graphicx}
\usepackage{booktabs}
\usepackage{algorithm}
\usepackage{multirow}
\usepackage[dvipsnames]{xcolor}
\usepackage{svg}
\usepackage{amsmath}
\usepackage{makecell}
\usepackage{pifont}
\newcommand{\xmark}{\ding{55}}
\usepackage{graphicx}
\usepackage{enumitem}
\DeclareMathOperator*{\argmax}{argmax}
\title{
Beyond Missing at Random: Incremental Tree Ensembles for Expanding Feature Sets
}

\author{Anonymous Author(s)}
\affiliation{\institution{Under Double-Blind Review} \country{Unknown}}
\email{anonymous@example.com}

\begin{document}

\begin{abstract}
Modern gradient boosting frameworks handle missing values through heuristic default-direction routing, treating absence as a nuisance rather than a signal. This approach fails under \textit{structured missingness} --- the common real-world pattern where features are systematically absent for identifiable subpopulations (e.g., diagnostic tests ordered only for critical patients, or financial metrics collected only for premium segments). Under structured missingness, the feature-absent and feature-present populations may exhibit fundamentally different predictive patterns, yet a single model must serve both through the same tree structure.

We introduce a population-aware learning framework that treats structured missingness as a modeling opportunity. Instead of routing around missing values, we explicitly partition training into subpopulation-specific phases: a base model specializes on instances lacking extended features, while a warm-started extended model incorporates newly available features through adaptive tree pruning. This two-phase design enables learning across evolving and distributed feature spaces without data sharing --- the base model artifact transfers knowledge between phases without exposing raw data.

Evaluation across 10 diverse datasets demonstrates that our approach matches or outperforms the standard two-model baseline --- which trains separate base and combined models with full data access --- while requiring no data sharing between phases (Wilcoxon $p = 0.002$, all 10 datasets positive). This proves that the two-phase approach is a viable alternative to standard training, and sometimes even superior, while additionally enabling deployment under data privacy constraints and evolving feature availability. All 10 datasets confirm successful model-level knowledge transfer without cross-phase data access.
\end{abstract}

\keywords{Structured Missingness, Population-Aware Learning, Gradient Boosting, Feature Evolution, Privacy-Preserving ML, Adaptive Tree Pruning, Model Transfer}

\maketitle

\section{Introduction}

Machine learning models are increasingly deployed in dynamic environments where the feature space evolves over time. In healthcare, patients may undergo new diagnostic tests during hospitalization, leading to heterogeneous feature availability across records \cite{rajkomar2018healthcare}. In finance, regulatory changes and behavioral metrics are introduced incrementally and collected only for certain user segments \cite{dixon2020finance}. These evolving schemas lead to structured, non-random missingness where new features are systematically absent from earlier records not due to noise, but as a result of how data is acquired. Importantly, these features do not update at uniform intervals, creating temporal asymmetry in data availability.

Tree-based models, particularly GBMs, remain the dominant choice in these domains due to their strong performance on tabular data \cite{rana2023comparative,langenberger2023hcp}, natural handling of mixed data types and missing values, and inherent explainability - a critical requirement in healthcare and finance applications.

Maintaining robust predictive models in such settings presents a major challenge. Organizations often resort to training multiple specialized models aligned with different feature subsets. However, this strategy introduces substantial maintenance burden, complicates deployment workflows, and risks version drift across models. Moreover, it is incompatible with real-world scenarios involving partial feature availability at inference time, particularly under privacy constraints such as GDPR where feature usage may vary at the user level.

A closely related challenge arises in cross-organizational collaboration, where multiple institutions hold complementary feature sets but cannot share raw data due to privacy regulations. For example, two hospitals may collect different diagnostic measurements for overlapping patient populations, or banks collaborating on anti-money laundering (AML) models may hold transaction features that cannot be exchanged. Our framework directly addresses this scenario: one institution trains a base model on its features and shares only the model artifact, while the other extends it with private features using warm-start and adaptive pruning. Our experimental design mirrors exactly this setup --- the base model is trained on one feature set and extended with additional features from a different subpopulation, with no data shared between phases.

In this work, we challenge the prevailing assumption that missing values should be handled heuristically within a single model. We show that when missingness is structured, explicit subpopulation specialization outperforms unified models across 10 diverse datasets (Wilcoxon $p = 0.002$). Our contributions are:
\begin{itemize}
\setlength\itemindent{0pt}
\setlength{\itemindent}{0pt}
\item \textbf{A modeling insight:} We provide empirical evidence that under structured missingness, the feature-absent population is systematically underserved by combined models that rely on heuristic missing-value routing. Training a dedicated base model for this subpopulation yields consistent AUC improvements across all 10 evaluated datasets. This finding holds regardless of dataset size (7K--568K), null rate (19--88\%), or domain.

\item \textbf{An adaptive knowledge transfer mechanism:} We develop a validation-driven tree pruning method that determines how much prior knowledge to retain when extending to new feature spaces. Bayesian optimization selects the pruning level, with a built-in $k=0$ fallback that guarantees the framework never underperforms training from scratch. Adaptive pruning outperforms fixed strategies in 8 of 10 datasets.

\item \textbf{A privacy-preserving collaboration paradigm:} Our two-phase architecture enables learning across distributed feature spaces without data sharing. Each experiment directly demonstrates this: the base model trains on one feature set and is extended with additional features from a different subpopulation, with no raw data exchanged. All 10 datasets validate successful model-level knowledge transfer, providing empirical evidence for deployment under GDPR, HIPAA, or institutional data-sharing constraints.
\end{itemize}

\section{Related Work}
Missing data are traditionally addressed either by discarding incomplete instances or by imputing missing values, but each strategy has drawbacks~\cite{chen2016xgboost}. Common strategies include MICE~\cite{vanbuuren2011mice}, which iteratively imputes values using regression under MAR assumptions, and tree-based methods like CART~\cite{loh2011cart} and MIA~\cite{twala2008mia}, which handle missingness via surrogate splits or by treating missing values as informative categories.

Modern gradient boosting frameworks such as XGBoost~\cite{chen2016xgboost}, LightGBM, and CatBoost include built-in mechanisms for missing values, but these are fundamentally heuristic. XGBoost learns a default split direction at each tree node: when a feature is absent during inference, the sample is routed left or right based on whichever direction minimized training loss. This approach treats missingness as ``just another value'' and optimizes the routing empirically. Critically, it does not model \textit{why} a value is missing, does not distinguish between random and structured missingness, and does not adapt the model architecture for subpopulations with different feature availability. When missingness is structured --- e.g., ICU lab tests ordered only for critical patients, or premium financial metrics collected only for high-value customers --- the missing-value population and the observed-value population may exhibit fundamentally different predictive patterns. A single model with heuristic default directions cannot specialize for these distinct subgroups.

Bayesian Additive Regression Trees (BART)~\cite{chipman2010bart} can be augmented with MIA to natively support incomplete data without imputation, but similarly assume that missingness is random or independent. This limitation motivates our approach, which explicitly models structured missingness by training separate specialized models for each subpopulation rather than relying on heuristic missing-value routing within a single model.
Recent research has begun addressing more complex streaming and evolving-feature scenarios, where both the data distribution and feature space can change over time. While traditional stream learners assume a fixed set of input features, real-world systems often experience column-wise feature drift with new sensors, logs, or sources introducing features that were previously unavailable.
Several ensemble based methods, such as Adaptive XGBoost (AXGB) and Aditer, adapt to row-wise concept drift by updating trees incrementally, often using sliding windows or tree replacement strategies. However, these models typically assume a static input space and cannot gracefully integrate new feature columns without full retraining. Other frameworks like Online learning from Capricious Data Streams (OCDS) \cite{he2019ocds} attempt to bridge feature sets by learning mappings between old and new features using generative models, but they still rely on conventional missing-value treatments and may struggle with structured missingness when the absence of data is systematic and correlated with the target. These limitations motivate our approach. Rather than retraining from scratch, we incrementally integrate new feature columns into an existing gradient boosting model, preserving previously learned structure and enabling adaptation to evolving schemas and feature sets - even as data distributions shift.

ORF3V is an online forest-based model designed to handle data streams with changing feature spaces~\cite{beyazit2019orf3v}. Each feature is tracked by its own mini-ensemble of shallow trees, allowing the model to add or remove features dynamically. Unlike our method, ORF3V maintains a separate model per feature and does not distinguish between base and extended phases. Our approach focuses on structured feature allocation and controlled model updates for deployment simplicity.
Adaptive XGBoost (AXGB)\cite{montiel2020axgb} updates an ensemble online by adding/replacing trees on mini-batches and triggering concept-drift detectors but assumes a fixed feature space, while Streaming Gradient Boosted Trees (SGBT)\cite{vaquet2022sgbt} embeds a tree-replacement mechanism for tracking distribution shifts in data streams with unchanged features; in contrast, our method tackles the complementary problem with a static population but growing columns, grafting additional features into the booster through extra boosting rounds and targeted pruning.
GBDT-IL is an incremental boosting method that updates the model over time by adding new trees and removing outdated ones~\cite{zhang2019gbdtil}. It selects useful features using the Fisher Score, but does not handle situations where some features are missing. In contrast, our method is designed for evolving feature sets and can adapt when only part of the features are available.
PUFE~\cite{he2021pufe} and Evolving Metric Learning (EML)~\cite{dong2021eml} both deal with changing feature spaces over time. PUFE fills missing overlaps using matrix completion and combines predictions from multiple models in an ensemble, while EML transforms vanished features into a compact form and integrates new ones using a learned distance metric in two stages. In contrast to these approaches, our method focuses on structured missingness in static tabular data where new features are added rather than vanishing. Our method extends a base GBM with new features in a second phase, preserving the original model without retraining or ensembling.
LIR-eGB. LIR-eGB~\cite{wang2023liregb} is designed for data streams where the distribution changes over time (concept drift). It improves a GBM by pruning weak trees that no longer help, based on a loss improvement ratio (LIR). The model switches between naive and statistical pruning, and adds new trees incrementally to adapt to the changing data. In contrast, our method focuses on static tabular data with structured missing features. We extend an existing GBM by adding new features in a second phase while preserving useful base trees, without retraining or handling concept drift.
\begin{table}[ht]
  \caption{Capability matrix (rows = method, columns = key capability)}
  \label{tab:method-comparison-transposed}
  \centering
  \resizebox{0.95\columnwidth}{!}{%
  \begin{tabular}{l c c c c c}
    \toprule
    & \makecell[c]{\textbf{Column-wise}\\\textbf{incremental}\\\textbf{learning}}
    & \makecell[c]{\textbf{Persistent}\\\textbf{feature}\\\textbf{support}}
    & \makecell[c]{\textbf{Handles}\\\textbf{structured}\\\textbf{missingness}}
    & \makecell[c]{\textbf{Single}\\\textbf{unified}\\\textbf{model}}
    & \makecell[c]{\textbf{Code}\\\textbf{Provided}}\\
    \midrule
    \textbf{AXGB} & \large\times & \large\times & \large\times &  \checkmark & \checkmark \\
    \textbf{SGBT} & \large\times & \large\times & \large\times & \checkmark & \checkmark \\
    \textbf{OCDS} & \checkmark & \checkmark & \checkmark & \large\times & \large\times \\
    \textbf{ORF3V} & \checkmark & \checkmark & \large\times & \large\times & \large\times\\
    \textbf{GBDT-IL} & \large\times & \checkmark & \large\times & \checkmark & \large\times \\
    \textbf{PUFE} & \checkmark & \large\times & \checkmark &  \checkmark & \large\times \\
    \textbf{LIR-eGB} & \large\times & \large\times & \large\times & \checkmark & \checkmark \\
    \textbf{EML}  & \checkmark & \large\times & \checkmark &  \checkmark & \large\times \\
    \textbf{Our}         &  \checkmark & \checkmark & \checkmark & \checkmark & \checkmark \\
    \bottomrule
  \end{tabular}%
  }
\end{table}
\section{Problem Formulation}
\label{sec:problem}

We formalise the setting in which column-wise incremental learning is required and specify the objectives that our method must meet.

\subsection{Data Model and Notation}

Let
\[
\mathcal D=\{(\mathbf{x}_i,y_i)\}_{i=1}^{N},
\qquad
y_i\in\{0,1\},
\]
be a binary-classification dataset collected over a time horizon
\(0<t_1<t_2<\dots<t_N\).
Each input vector \(\mathbf{x}_i\in\mathbb R^{p_B+p_E}\) is partitioned into
\[
\mathbf{x}_i=\bigl[\;\mathbf{b}_i\;\|\;\mathbf{e}_i\bigr],
\qquad
\mathbf{b}_i\in\mathbb R^{p_B},\;
\mathbf{e}_i\in\mathbb R^{p_E},
\]
where \(\mathbf{b}_i\) is the base-feature block that is always observed, and
\(\mathbf{e}_i\) is the extended-feature block, present only for instances that arrive
after a feature-release time \(t_E\) (or when acquisition conditions are met).

A binary missingness mask
\(\mathbf{m}_i\in\{0,1\}^{p_E}\) is defined by
\(m_{ij}=1\) if the \(j\)-th extended feature of instance \(i\) is observed and
\(m_{ij}=0\) otherwise.
We assume structured (NMAR) missingness:
\[
P(\mathbf{m}_i \mid \mathbf{b}_i,\mathbf{e}_i,y_i)
\;\neq\;
P(\mathbf{m}_i \mid \mathbf{e}_i,y_i),
\]
i.e.\ the probability that a value is missing may depend on observed base features (
In healthcare, advanced laboratory tests are ordered only when certain vitals exceed thresholds; in finance, premium behavioural metrics are gathered only for high-value customers).

\subsection{Predictive Objective}

We seek a single probabilistic classifier
\[
f_\theta:\;
  \mathbb R^{p_B+p_E}\times\{0,1\}^{p_E}
  \;\longrightarrow\;[0,1],
\]
parameterised by \(\theta\), that satisfies:

\begin{enumerate}[label=(\alph*)]
\item \textbf{Validity under partial information.}\;
For any instance \(i\), inference uses only the available features:
\[
\hat y_i
  = f_\theta\!\Bigl(
      \bigl[\mathbf{b}_i \,\|\, \mathbf{m}_i\odot\mathbf{e}_i\bigr],
      \mathbf{m}_i
    \Bigr),
\]
where "\(\odot\)" denotes element-wise multiplication (missing entries are zeroed).

\item \textbf{Minimisation of predictive error.}\;
The optimisation target is the binary log-loss evaluated on a validation set:
\[
\mathcal L_{\mathrm{pred}}(\theta)=
\frac{1}{|\mathcal D_{\mathrm{val}}|}
\sum_{(\mathbf{x}_i,y_i)\in\mathcal D_{\mathrm{val}}}
\Bigl[
  -y_i\log\!\bigl(\hat y_i\bigr)
  -(1-y_i)\log\!\bigl(1-\hat y_i\bigr)
\Bigr],
\]
with \(\hat y_i\) defined as above.
\end{enumerate}

\subsection{Incremental-Learning Constraint}

Model training proceeds in two phases.  Let \(\theta^{(0)}\) denote the parameters of
a base model trained solely on \(\{\mathbf{b}_i\}\).  Extending this model
yields
\[
\theta^{\star}
=
\arg\!\min_{\theta}
\;
\underbrace{\mathcal L_{\mathrm{pred}}(\theta)}_{\text{predictive performance}}
\;+\;
\lambda\,\underbrace{C_{\mathrm{inc}}\!\bigl(\theta\mid\theta^{(0)}\bigr)}_{\text{incremental cost}},
\tag{1}\label{eq:inc-opt}
\]
where \(C_{\mathrm{inc}}(\theta\mid\theta^{(0)})\) measures the additional computation
required to transform \(\theta^{(0)}\) into the extended-feature model \(\theta\),
and \(\lambda>0\) trades off predictive gain against update cost.

\paragraph{Deployment guarantees.}
For optimisation~\eqref{eq:inc-opt} to be acceptable in production we require:

\begin{enumerate}[label=(\Roman*)]
\item No degradation on legacy traffic:
\[
\operatorname{AUC}\!\bigl(f_{\theta^{\star}};\{\mathbf{m}_i=\mathbf 0\}\bigr)
\;\ge\;
\operatorname{AUC}\!\bigl(f_{\theta^{(0)}};\{\mathbf{m}_i=\mathbf 0\}\bigr).
\]

\item \emph{Monotonic benefit on enriched traffic:}
for the subset with any \(m_{ij}=1\),
\[
\operatorname{AUC}\!\bigl(f_{\theta^{\star}}\bigr)
\;\ge\;
\operatorname{AUC}\!\bigl(f_{\theta^{(0)}}\bigr).
\]
\end{enumerate}

Given \((\mathcal D,\mathbf m)\) and trade-off \(\lambda\), our goal is to design an
\textbf{incremental training procedure} that solves~\eqref{eq:inc-opt} whenever new
columns appear, thereby delivering a single deployable model
\(f_{\theta^{\star}}\) that: (i) attains near-optimal AUC on both basic and
extended cohorts, (ii) controls incremental cost through adaptive pruning and
continued boosting, and (iii) gracefully handles heterogeneous features
availability.

\section{METHODS}
\label{sec:method}

Our approach implements a population-aware gradient boosting framework for learning across evolving and distributed feature spaces without data sharing, as illustrated in Figure~\ref{fig:framework}. Rather than training a single model that heuristically handles missing values, we explicitly partition the learning process into subpopulation-specific phases. The framework consists of two main components: Base Model Training (specializing on the feature-absent population) and Extended Model Training (incorporating newly available features through knowledge transfer). Bayesian optimization determines the optimal feature allocation between phases.

\subsection{Phase 1: Base Model Training}

First, we construct a base model using a gradient boosting machine (GBM) trained exclusively on commonly available features. This foundation model undergoes thorough hyperparameter tuning through Bayesian optimization to maximize performance on the initial feature set. The optimized hyperparameters include learning rate, tree depth, regularization parameters, and other model-specific configurations to ensure optimal validation set performance and generalizability.

The underlying predictive model utilizes gradient boosting machines (GBMs), an ensemble learning method that builds a series of decision trees sequentially, with each tree correcting errors made by previous trees. The base model is formulated as:

\begin{equation}
F_{base}(x) = \sum_{i=1}^{n} \alpha_i h_i(x_{F_{base}})
\end{equation}

Where $h_i$ are decision trees trained on the base features $F_{base}$, and $\alpha_i$ are tree weights.

\subsection{Phase 2: Extended Model Training}

The second phase develops an extended model that incorporates newly available features present in only a subset of instances, corresponding to Phase 2 in Figure~\ref{fig:framework}. Rather than training from scratch, we employ a transfer learning approach with three key mechanisms:

\begin{enumerate}
\item \textbf{Warm-start initialization} from the base model, preserving knowledge gained from widely available features

\item \textbf{Adaptive tree pruning} to determine the optimal subset of base trees to retain, strategically balancing prior knowledge with flexibility for new patterns

\item \textbf{Continued training} that allows the model to learn relationships involving the extended feature set
\end{enumerate}

The pruning level determining exactly how many trees $k$ to preserve from the base model is itself optimized through Bayesian optimization against a validation set. This process evaluates different pruning thresholds to identify the configuration that best transfers relevant knowledge while avoiding concept drift between the feature spaces.

The extended model builds upon this foundation through transfer learning:

\begin{equation}
F_{ext}(x) = \sum_{i=1}^{k} \alpha_i h_i(x_{F_{base}}) + \sum_{j=1}^{m} \beta_j g_j(x_{F_{base} \cup F_{ext}})
\end{equation}

Where $k \leq n$ is the optimal number of base trees to retain (determined through pruning), and $g_j$ are additional trees trained on the combined feature space.

\subsection{Feature Allocation using Bayesian Optimization}

We optimize the allocation of features between base and extended sets by running both Phase 1 and Phase 2 with different feature allocation configurations using Bayesian optimization. This optimization process evaluates different partitioning of features between base and extended sets to determine the optimal feature allocation strategy.

The feature allocation optimization aims to minimize the objective function:

\begin{equation}
\mathcal{L} = \text{AUC}_{\text{baseline}} - \text{AUC}_{\text{two-phase}}
\end{equation}

Where $\text{AUC}_{\text{baseline}}$ is the validation area under the ROC curve for a model trained with all features at once, and $\text{AUC}_{\text{two-phase}}$ is the validation performance of our two-phase approach. The Bayesian optimization process systematically explores different feature allocation strategies to identify the configuration that maximizes the performance advantage of our two-phase method over the baseline approach.

In real-world scenarios, feature allocation is subject to the availability of data sources at training time, making this optimization crucial for determining the most effective partitioning strategy given the constraints of data accessibility and acquisition timing.

\subsection{Inference}

The inference process, depicted in the lower section of Figure~\ref{fig:framework}, demonstrates how new data is routed to either the base or extended model depending on feature availability. During inference, the system automatically determines which model configuration to use based on the available features for each instance:

\begin{itemize}
\item \textbf{Base model inference}: Applied when only base features are available
\item \textbf{Extended model inference}: Applied when both base and extended features are present
\item \textbf{Adaptive routing}: The system seamlessly switches between models based on real-time feature availability
\end{itemize}

This comprehensive approach enables ongoing model enhancement as new features emerge, without requiring complete retraining or sacrificing performance on instances with limited feature availability.

\subsection{Theoretical Analysis - Performance Guarantee}

Our two-phase incremental learning framework provides a fundamental performance guarantee: the extended model will not perform worse than training a new model from scratch. This guarantee is ensured by our adaptive pruning mechanism in Algorithm 1.

The algorithm evaluates pruning levels $k \in \{0, 1, 2, \ldots, |T_{\text{base}}|\}$ on validation data and selects the optimal $k$ that maximizes performance. Crucially, $k=0$ corresponds to using no base trees and training a completely new model from scratch on the extended dataset.

This theoretical foundation ensures that our incremental approach is always a safe choice - it will either successfully transfer useful knowledge from the base model (when $k > 0$ is selected) or automatically fall back to the baseline performance of training from scratch (when $k = 0$ is selected).

\begin{algorithm}[t]
\caption{Adaptive Tree Pruning for Incremental Learning in Gradient Boosting}
\begin{algorithmic}[1]
\Require Base model $M_{base}$ with trees $T_{base} = \{t_1,...,t_n\}$
\Require New data batch $D_{new}$ with features $F$
\Ensure Updated model $M_{updated}$
\Function{FindOptimalPruning}{$M_{base}$, $D_{new}$}
    \State $D_{train}, D_{val} \gets$ SplitData($D_{new}$)
    \State $scores \gets [\ ]$

    \For{$k = 0$ \textbf{to} $|T_{base}|$}
        \State $M_k \gets$ Initialize with first $k$ trees (empty if $k=0$)
        \State $scores[k] \gets$ Evaluate($M_k$, $D_{val}$)
    \EndFor

    \State \Return $\argmax_k(scores)$
\EndFunction
\State $k_{opt} \gets$ \Call{FindOptimalPruning}{$M_{base}$, $D_{new}$}
\If{$k_{opt} = 0$}
    \State $M_{updated} \gets$ Train new model from scratch on $D_{new}$
\Else
    \State $M_{updated} \gets$ Initialize with trees $\{t_1,...,t_{k_{opt}}\}$
    \State Add new trees to $M_{updated}$ by training on $D_{new}$

\EndIf
\State \Return $M_{updated}$
\end{algorithmic}
\end{algorithm}


\begin{figure}
  \centering
  \includegraphics[width=0.5\textwidth]{1inc.png}
  \caption{A two-phase framework: (1) Base Model Training builds a foundation using common features; (2) Extended Model Training enhances this with newly available features using adaptive pruning and drift detection. During inference, data routes to either model based on feature availability.}

  \label{fig:framework}
\end{figure}

\section{Experiments}
This section evaluates our population-aware learning framework across diverse real-world datasets. We demonstrate how explicit subpopulation specialization, combined with adaptive tree pruning, consistently outperforms a single combined model that relies on heuristic missing-value handling. We conduct ablation studies to validate the contribution of the adaptive pruning mechanism and present a systematic analysis of what drives the observed improvements.

All Bayesian optimization procedures use Optuna\cite{optuna} with 30 trials per model (50 for select sensitivity experiments). Training and validation employ 3-fold stratified cross-validation to ensure robust performance estimates and prevent overfitting, particularly for the imbalanced class distributions in our datasets.
\subsection{Datasets}
We evaluate our methodology on 10 diverse binary classification datasets spanning healthcare, finance, meteorology, and behavioral prediction domains (Table~\ref{tab:dataset-summary}). The datasets vary significantly in size (7K to 568K rows), feature dimensionality (11 to 183 columns), class balance (1.9\% to 77.4\% positive), and missingness patterns (19\% to 88\% null rate among extended features).

Six datasets contain naturally occurring structured missingness: Weather and WeatherAUS (meteorological measurements missing due to station equipment availability), DiabetesRecord (clinical measurements absent for certain patient subgroups), HRAnalytics (demographic fields with selective disclosure), WIDS (ICU hour-1 lab values ordered only for critical patients, with 83--92\% null rates), and ClientRecordAug (service-dependent features absent for non-subscribers). Four datasets use controlled null injection to simulate feature evolution scenarios: BankLoanSta, ClientRecordV2, MovieAugV2, and FlightDelay, where we introduce independent random missingness at rates of 20--50\% into selected features while preserving the original data relationships.

This combination of natural and augmented missingness allows us to evaluate our framework under both realistic deployment conditions and controlled experimental settings. Each dataset underwent domain-specific preprocessing including removal of identifier columns, target encoding, and feature type standardization.

Figure~\ref{fig:experimental_workflow} illustrates our experimental workflow. The process partitions each dataset, separates base and extended features according to missingness patterns, and implements the two-phase training procedure with adaptive pruning.

\begin{figure}
  \centering
  \includegraphics[width=0.48\textwidth]{incremental-tree-diagram_2.png}
  \caption{Experimental workflow for incremental feature-based learning. The process begins with splitting an original dataset into training and test sets (top), followed by separating base and extended features. The base model is trained on commonly available features, then extended through adaptive pruning to incorporate newly available features. Performance evaluation compares AUC between our method and a baseline}

  \label{fig:experimental_workflow}
\end{figure}

\begin{table}[h]
\centering
  \label{tab:dataset-summary}
 \caption{
    Datasets used in experiments. Label\% shows positive class proportion. Null\% is the average missingness of the extended features used in each experiment.
  }\resizebox{0.95\columnwidth}{!}{%
  \begin{tabular}{l r r r r}
    \toprule
    \textbf{Dataset} & \textbf{Rows} & \textbf{Columns} & \textbf{Label\%} & \textbf{Null (\%)} \\
    \midrule
    FlightDelay        & 567,630  & 14  & 17.4 & 20.0 \\
    WeatherAUS         & 142,193  & 20  & 22.4 & 40.3 \\
    MovieAugV2         & 112,466  & 11  & 1.9  & 50.0 \\
    DiabetesRecord     & 101,766  & 45  & 46.1 & 66.2 \\
    BankLoanSta        & 100,000  & 17  & 77.4 & 19.1 \\
    WIDS               & 91,713   & 183 & 8.6  & 87.7 \\
    Weather            & 25,000   & 16  & 22.4 & 49.3 \\
    HRAnalytics        & 19,158   & 14  & 24.9 & 23.5 \\
    ClientRecordAug    & 7,043    & 28  & 67.0 & 55.0 \\
    ClientRecordV2     & 7,043    & 28  & 67.0 & 47.5 \\
  \bottomrule
  \end{tabular}%
  }


\end{table}


\subsection{Experimental Set}
For each dataset, we implemented a five step experimental process to evaluate our incremental feature based learning approach:

\begin{enumerate}
\item \textbf{Preprocessing and Feature Optimization}: We performed an optimization routine to determine the most effective allocation between base and extended feature sets, balancing predictive power against missing value patterns.

\item \textbf{Base Model Construction}: We trained an XGBoost model using only the optimized base feature set. Hyperparameters were tuned using an  bayesian optimization, resulting in optimal configurations for each dataset.

\item \textbf{Extended Model Development}: We constructed an extended model by:
\begin{itemize}
    \item Using the base model as initialization
    \item Applying adaptive tree pruning to selectively retain beneficial trees from the base model
    \item Determining the optimal pruning level through validation performance
    \item Continuing training with instances containing extended features
\end{itemize}

\item \textbf{Test Group Stratification}: We divided the test dataset into two groups:
\begin{itemize}
    \item Base Group: Instances without extended features
    \item Extended Group: Instances with non-null extended features
\end{itemize}

\item \textbf{Performance Evaluation}: We calculated AUC for each test group, comparing our approach against a standard two-model baseline: a base model trained on base features only, and a combined model trained on all features with full data access. Both approaches serve the same two test populations; the difference is that our method can operate with data separation between phases while the baseline cannot.
\end{enumerate}

As described in Section 4.2, the tree pruning strategy was essential for effective model combination, allowing us to identify and retain only the most relevant decision trees when transitioning from the base to the extended model.

\textbf{Ablation Study}: To isolate the contribution of our adaptive tree pruning mechanism, we conducted an ablation study comparing models with and without the pruning component. This analysis evaluated the effectiveness of selectively retaining base trees versus preserving the entire base model when incorporating extended features.

\section{Results}

Our population-aware framework demonstrates consistent improvements across a diverse set of 10 real-world datasets. Table~\ref{tab:merged_auc_results} presents a detailed comparison. Both approaches serve two test populations: instances without extended features (base population) and instances with extended features (extended population). The key difference is:

\begin{itemize}
\item \textbf{Standard approach (baseline):} Trains a base model using only base features, and a separate combined model using all features. The base model serves the base population; the combined model serves the extended population. This requires simultaneous access to all data and features.
\item \textbf{Our approach:} Trains a base model on all rows with extended features set to NaN, then creates an extended model by warm-starting from the base model and training on the extended subpopulation with adaptive tree pruning. The two phases can operate independently on separate data, enabling privacy-constrained deployment.
\end{itemize}

\subsection{Comprehensive Performance Analysis}

Across all 10 datasets, our framework achieves positive total improvement, demonstrating that the two-phase incremental approach --- which does not require simultaneous access to all data --- matches or outperforms the standard approach that does. The strongest gains appear in BankLoanSta (+0.0743), Weather (+0.0097), FlightDelay (+0.0046), and ClientRecordV2 (+0.0046).

A key finding is that our base model consistently outperforms the baseline's base model on the no-extended test population (positive base diff in all 10 datasets). This occurs because our base model trains on \textit{all} rows with extended features masked as NaN, giving it a larger training set than the baseline's base model which uses only base features. The extended features' absence pattern itself becomes an implicit signal that aids specialization.

For the extended population, results are more nuanced. In 6 out of 10 datasets, the extended model matches or exceeds the combined baseline on extended-feature rows. In the remaining 4 datasets (Weather, HRAnalytics, ClientRecordAug, WeatherAUS), the extended model shows marginal deficits (less than 0.002 AUC), which are more than compensated by the base model's gains on the no-extended population.

\textbf{Statistical significance.} All 10 datasets show positive total improvement. A Wilcoxon signed-rank test confirms that the improvements are statistically significant ($W = 0$, $p = 0.002$), indicating that the probability of all 10 datasets showing positive improvement by chance is less than 0.2\%. The median improvement across datasets is +0.0025 AUC. We note that a parametric $t$-test does not reach significance ($p = 0.20$) due to the high variance introduced by the BankLoanSta outlier (+0.0743), underscoring the importance of non-parametric testing when dataset-level effects vary in magnitude.

\textbf{Robustness across pruning strategies.} As a proxy for result stability, we examine the spread of improvements across the three pruning strategies tested in our ablation study (Table~\ref{tab:pruning}). While the adaptive (Optuna) strategy achieves the best improvement for each dataset, performance varies across strategies: the mean range between best and worst pruning mode is 0.004 AUC. Two datasets (BankLoanSta and ClientRecordV2) show positive improvement under all three pruning strategies, indicating robust gains regardless of pruning choice. For the remaining datasets, at least one strategy achieves positive improvement, while others may yield marginal deficits --- confirming that pruning strategy selection matters and justifying the adaptive optimization approach.

\begin{table*}[!t]
\caption{AUROC Results Across All Datasets. Both approaches use two models evaluated on separate test populations. \textit{Baseline}: trains a base model on base features only, and a combined model on all features --- requires access to all data. \textit{Our approach}: trains a base model on base features (with extended features as NaN), then warm-starts an extended model using adaptive pruning --- phases can operate on separate data. Total Improvement is the sum of AUC gains across both populations.}
\label{tab:merged_auc_results}
\centering
\small
\begin{tabular}{l r r c c c c c}
\toprule
\textbf{Dataset} & \textbf{Rows} & \textbf{Null\%} & \textbf{Our (base)} & \textbf{Baseline (base)} & \textbf{Our (ext)} & \textbf{Baseline (ext)} & \textbf{Total Improvement} \\
\midrule
BankLoanSta & 100K & 19.1 & 0.6952 & 0.6495 & 0.7817 & 0.7531 & \textcolor{ForestGreen}{+0.0743} \\
Weather & 25K & 49.3 & 0.8432 & 0.8314 & 0.8448 & 0.8469 & \textcolor{ForestGreen}{+0.0097} \\
DiabetesRecord & 102K & 66.2 & 0.6734 & 0.6727 & 0.6682 & 0.6663 & \textcolor{ForestGreen}{+0.0026} \\
HRAnalytics & 19K & 23.5 & 0.7735 & 0.7727 & 0.8102 & 0.8108 & \textcolor{ForestGreen}{+0.0002} \\
ClientRecordAug & 7K & 55.0 & 0.9612 & 0.9605 & 0.9608 & 0.9612 & \textcolor{ForestGreen}{+0.0004} \\
ClientRecordV2 & 7K & 47.5 & 0.9235 & 0.9198 & 0.9186 & 0.9177 & \textcolor{ForestGreen}{+0.0046} \\
MovieAugV2 & 112K & 50.0 & 0.9187 & 0.9185 & 0.8726 & 0.8726 & \textcolor{ForestGreen}{+0.0001} \\
WeatherAUS & 142K & 40.3 & 0.8865 & 0.8861 & 0.9024 & 0.9025 & \textcolor{ForestGreen}{+0.0003} \\
WIDS & 92K & 87.7 & 0.8997 & 0.8991 & 0.8967 & 0.8950 & \textcolor{ForestGreen}{+0.0024} \\
FlightDelay & 568K & 20.0 & 0.7478 & 0.7450 & 0.7188 & 0.7170 & \textcolor{ForestGreen}{+0.0046} \\
\bottomrule
\end{tabular}
\end{table*}

\subsection{Ablation Study - Impact of Adaptive Tree Pruning}

We conducted an ablation study comparing three pruning strategies across all 10 datasets (Table~\ref{tab:pruning}): (1) \textit{Optuna} --- Bayesian optimization selects whether to warm-start and how many trees to retain; (2) \textit{No Pruning} --- always warm-starts from the full base model; (3) \textit{Fixed 50\%} --- always retains exactly half the base trees.

Optuna-based adaptive pruning achieves the best objective in 8 out of 10 datasets, confirming that automatic pruning threshold selection outperforms fixed strategies. No Pruning wins in 2 datasets (HRAnalytics and ClientRecordAug), both of which are smaller datasets where the full base model transfers well. Fixed 50\% pruning never achieves the best result and produces positive objectives (combined model wins) in 6 out of 10 datasets, indicating that aggressive fixed pruning often discards useful knowledge.

Critically, the Optuna strategy includes the option of pruning all base trees ($k=0$), which is equivalent to training from scratch. This built-in fallback ensures that the incremental approach never underperforms a full retrain, providing a safety guarantee for production deployment.

\begin{table}[h]
\centering
\caption{Ablation Study: Pruning Mode Comparison. Objective = (combined$_\text{ext}$ -- extended AUC) + (combined$_\text{no}$ -- base AUC). Negative values indicate incremental learning outperforms the combined baseline. Bold marks the best mode per dataset.}
\label{tab:pruning}
\resizebox{0.95\columnwidth}{!}{%
\begin{tabular}{l r r r l r}
\toprule
\textbf{Dataset} & \textbf{Optuna} & \textbf{No Pruning} & \textbf{Fixed 50\%} & \textbf{Best} & \textbf{Ext AUC} \\
\midrule
BankLoanSta & \textbf{--0.0743} & --- & --- & Optuna & 0.7817 \\
Weather & \textbf{--0.0097} & +0.0081 & +0.0074 & Optuna & 0.8448 \\
DiabetesRecord & \textbf{--0.0026} & +0.0006 & --0.0011 & Optuna & 0.6682 \\
HRAnalytics & --0.0002 & \textbf{--0.0009} & +0.0011 & No Pruning & 0.8109 \\
ClientRecordAug & --0.0004 & \textbf{--0.0005} & +0.0006 & No Pruning & 0.9610 \\
ClientRecordV2 & \textbf{--0.0046} & --0.0038 & --0.0034 & Optuna & 0.9186 \\
MovieAugV2 & \textbf{--0.0040} & +0.0009 & +0.0052 & Optuna & 0.8739 \\
WeatherAUS & \textbf{--0.0003} & --0.0003 & +0.0016 & Optuna & 0.9024 \\
WIDS & \textbf{--0.0045} & --0.0024 & +0.0003 & Optuna & 0.8955 \\
FlightDelay & \textbf{--0.0046} & --0.0043 & +0.0023 & Optuna & 0.7188 \\
\midrule
\multicolumn{4}{l}{\textbf{Wins}} & \multicolumn{2}{l}{Optuna: 8, No Pruning: 2, Fixed 50\%: 0} \\
\bottomrule
\end{tabular}%
}
\end{table}

\subsection{Analysis: What Drives Incremental Learning Success?}

To understand when and why our incremental approach outperforms the combined baseline, we conducted a systematic correlation analysis across all 153 experiment configurations. We examined whether measurable dataset or experiment properties --- including null rates, dataset size, label distribution, feature importance ratios, and population shift between extended and non-extended subgroups --- predict the objective.

Our analysis reveals that no single dataset-level property strongly predicts the improvement magnitude ($R^2 = 0.019$ for a multivariate regression on all measured properties). Feature importance of extended features in the combined model shows no correlation with the objective ($r = -0.035$, $p = 0.68$). Similarly, covariate shift between the extended and non-extended populations, measured via Kolmogorov-Smirnov statistics and label rate differences, shows no significant relationship.

A decomposition of the objective into its two components provides an important insight. The base model outperforms the combined baseline on the no-extended population in 59.5\% of experiments, while the extended model outperforms the combined baseline on the extended population in only 31.4\% of experiments. This indicates that the primary source of improvement is the \textit{specialization of the base model} on instances lacking extended features --- a population that the combined model must share capacity across both subgroups to serve.

These findings suggest that the benefit of our framework is architectural rather than data-dependent: the two-phase design inherently enables population-specific specialization that a single combined model cannot achieve, regardless of the specific dataset characteristics.

\section{Discussion}

Our experimental results demonstrate the effectiveness of incremental feature based learning across diverse datasets. However, several important considerations merit further discussion regarding the method's applicability and limitations.

\subsection{Comparison with Existing Approaches}

Our experimental evaluation compared the proposed incremental feature based learning approach against baseline implementations from the River Python package\cite{montiel2021river}, specifically using a two phase ARFClassifier and AdaBoostClassifier. However, The underlying algorithmic differences between our Gradient Boosting Machine approach and the Random Forest/AdaBoost baselines introduce factors to isolate the specific contribution of our incremental learning strategy. Experimental results, and comparative analyses can be found in the supplementary materials.


In Addition, we emulate PUFE's ensemble approach by training two separate GBM models on disjoint feature subset. One on the base features and one on the extended features. During inference, we use whichever model's features are available. For instances where both feature sets are present, we average the predictions from both models. This approximation captures PUFE's late fusion strategy while avoiding the complexity of matrix completion.

Our experimental results (included in the supplementary materials) show that the PUFE-style ensemble yields slightly lower AUC on datasets with extended features compared to a single GBM trained on the full feature set. This may be attributed to PUFE's original design being more tailored to scenarios involving vanishing features, rather than expanding feature spaces. Thus, while PUFE is effective in handling fragmented or missing views, it may be less suitable when the feature set grows over time and unified training is feasible.

As shown in Table 3, We determined that comparing our approach against a model built from scratch on all features provides the most appropriate benchmark for evaluating our method. This comparison framework maintains consistency in the underlying algorithms and implementation details, with the only difference being the application of our two phase mechanism. This controlled experimental design allows us to isolate and quantify the specific impact of our incremental learning strategy across multiple diverse datasets and feature configurations.

\subsection{Why Two-Phase Matches or Outperforms the Standard Approach}

Our results demonstrate that the two-phase approach --- which can operate without simultaneous access to all data --- achieves performance comparable to or better than the standard approach that requires full data access. Two mechanisms drive this:

First, our base model trains on all rows with extended features as NaN, whereas the baseline's base model trains only on base features. This gives our base model access to more training instances and allows XGBoost's learned default directions to implicitly encode the missingness pattern as a signal. The result is a stronger base model that outperforms the baseline on the no-extended population in all 10 datasets.

Second, the warm-started extended model inherits knowledge from the base model and refines it with extended features, rather than learning from scratch. Adaptive pruning determines how much base knowledge to retain, preventing negative transfer while preserving useful patterns. This achieves comparable performance to the baseline's combined model on the extended population, while requiring only the model artifact (not the original training data) from Phase 1.

The broader implication is significant for deployment: organizations can achieve standard-approach performance levels while maintaining data separation between phases --- enabling privacy-constrained collaboration and incremental model updates as new features become available over time.

\subsection{Framework Contributions}

Despite the performance advantages, we emphasize that the main contribution of our work is the framework itself, which offers improved maintainability. The modular design allows for easier updates, monitoring, and extension compared to existing incremental learning solutions. Our framework also incorporates a critical rollback mechanism by concept, where in cases of significant data drift, it will rebuild the model from scratch rather than using the base model as input. This is accomplished by pruning all trees in the base model, effectively ensuring that outdated or irrelevant patterns don't negatively impact performance when data characteristics change substantially. This automatic adaptation mechanism provides an important safeguard for maintaining model reliability in dynamic environments.

\subsection{Application to Privacy-Constrained Model Collaboration}

Beyond feature evolution within a single organization, our experimental setup directly demonstrates a privacy-constrained collaboration scenario. In each experiment, the base model is trained on one feature set and the extended model incorporates additional features from a different subpopulation --- no raw data is shared between the two phases, only the base model artifact. This mirrors a real-world deployment where Institution A trains a base model on its available features and shares the model file, while Institution B extends it with private features that Institution A never accesses.

All 10 datasets in our evaluation validate this capability: the two-phase approach consistently matches or outperforms a combined model that would require access to all features simultaneously. This provides empirical evidence that model-level knowledge transfer can substitute for data-level sharing in cross-institutional settings such as healthcare (hospitals with different diagnostic panels) or finance (banks with complementary transaction features for AML detection).

We note that sharing model artifacts does not constitute formal privacy preservation --- gradient boosting models can be susceptible to model inversion attacks~\cite{fredrikson2015model}. Production deployments under strict regulatory constraints (GDPR, HIPAA) would benefit from additional safeguards such as differential privacy or secure aggregation. Nonetheless, the framework reduces the attack surface significantly compared to raw data sharing, and the model-sharing paradigm represents a practical middle ground between full data exchange and complete isolation.

\subsection{Generalizability and Limitations}

While all datasets in our experiments focus on binary classification, the methodology naturally extends to regression and multi-class tasks. For regression, the same adaptive pruning mechanism applies with appropriate loss functions. For multi-class classification, the approach can be implemented via one-vs-rest decomposition or native multi-class boosting with minimal modification.

A limitation of the current framework is that the optimal feature allocation between base and extended phases must be determined empirically through Bayesian optimization. Our analysis shows that no measurable dataset property predicts which allocation will succeed, suggesting that this search step cannot be easily shortcut. Additionally, run-to-run variance from the hyperparameter search process can be substantial, particularly for smaller datasets, requiring multiple runs for stable estimates.

\subsection{Future Work}

Several directions merit further investigation. First, extending the framework to support \textit{multiple sequential feature expansions} (Phase 3, Phase 4, etc.) would better reflect real-world scenarios where features arrive in multiple waves over time. Second, developing theoretical guarantees for when warm-starting provably outperforms training from scratch would reduce reliance on empirical validation. Third, exploring whether information-theoretic properties of the extended features (e.g., conditional mutual information with the target given base features) can predict incremental learning success would enable more principled feature allocation without exhaustive search. Finally, integrating our approach with privacy-preserving techniques such as differential privacy or secure aggregation could formalize the model-sharing paradigm described above.

\section{Conclusion}

We presented a population-aware learning framework for Gradient Boosting Machines that addresses the challenge of learning across evolving and distributed feature spaces without data sharing. Our approach treats structured missingness as a modeling opportunity: rather than heuristically routing around missing values, we train specialized models for subpopulations defined by feature availability, connected through adaptive tree pruning and knowledge transfer.

Evaluation across 10 diverse datasets from healthcare, finance, and meteorology domains demonstrates that our approach consistently matches or outperforms the standard approach that requires simultaneous access to all data, with total AUC gains ranging from +0.0001 to +0.0743 (Wilcoxon $p = 0.002$). This result is significant because our method achieves comparable performance while enabling data separation between training phases --- a property the standard approach cannot offer.

The adaptive pruning mechanism achieves the best results in 8 out of 10 datasets, with a built-in safety guarantee that ensures the framework never underperforms training from scratch. Our systematic analysis of 153 experiment configurations shows that the framework's benefit is architectural, arising from the two-phase design itself rather than from specific dataset characteristics.

The practical implications are twofold. For \textit{evolving feature spaces}, the framework enables incremental model updates as new features become available, without requiring access to historical training data or full retraining. For \textit{privacy-constrained collaboration}, institutions can share model artifacts rather than raw data --- all 10 datasets validate that this model-level knowledge transfer achieves consistent improvements without cross-phase data access. Together, these capabilities address a growing need in healthcare, finance, and other regulated domains where both data evolution and data privacy are first-order concerns.

\section*{Supplementary Materials and Code Availability}

To support reproducibility and facilitate further research, we provide supplementary materials that include:

\begin{itemize}
  \item The full implementation of our two-phase incremental GBM framework.
  \item Experiments code, logs and hyperparameter configurations for all datasets.
  \item Sensitivity analysis.
  \item Ablation study.
  \item Link: https://anonymous.4open.science/r/Learning-to-Grow-Incremental-Trees-over-Expanding-Features-818A/
\end{itemize}

\section*{Generative AI Usage Disclosure}

In accordance with ACM's Policy on Authorship and the guidelines for disclosure of generative AI usage, the authors hereby declare the following:


\noindent\textbullet\ \textbf{Language editing.} We used OpenAI ChatGPT and Anthropic Claude solely to improve spelling, grammar, punctuation, clarity, and overall readability of the manuscript. All suggested edits were critically reviewed, revised as necessary, and confirmed by the authors to ensure that the text remains an accurate and faithful representation of our original intellectual contributions.\\[0.5em]
\noindent\textbullet\ \textbf{Code development.} We employed OpenAI ChatGPT and Anthropic Claude to help scaffold boilerplate and to accelerate routine implementation tasks. Every line of code generated or suggested by these tools was examined, adapted, and tested by the authors to guarantee that it performs exactly as intended.\\[0.5em]

\noindent The authors take full responsibility for the veracity and correctness of all material in this Work, and no generative AI system is listed as an author.


\bibliographystyle{ACM-Reference-Format}
\bibliography{references}

\end{document}
